
\noindent
Zur weiteren Analyse der Leistungsübertragung wird auch bei der realen Stromquelle
die Leistung normiert dargestellt. Dadurch lassen sich die Ergebnisse unabhängig von
den konkreten Werten des Quellstroms und des Innenleitwerts verallgemeinern. Die
relative Leistung wird definiert als:
\begin{equation}
p = \frac{P}{P_0}
\end{equation}

\noindent
Als dimensionslose Größe wird erneut das Verhältnis zwischen Lastleitwert und Innen-
leitwert eingeführt:
\begin{equation}
r = \frac{G_L}{G_I}
\end{equation}

\noindent
Durch Einsetzen der Leistungsfunktion $P(G_L)$ aus Gleichung~\eqref{eq:PL_GL} sowie
der maximalen Leistung $P_0$ aus Gleichung~\eqref{eq:P0_stromquelle} in die Definition
der relativen Leistung ergibt sich:
\begin{equation}
p(r) = \frac{4r}{(1+r)^2}
\end{equation}
\begin{figure}[H]
    \centering
    \includegraphics[width=1\textwidth]{2_2_3_p_r_stromquelle.png}
    \caption{Leistung im Verhältnis der Widerstände (Stromquelle)}
    \label{fig:Leistung im Verhältnis der Widerstände (Stromquelle)}
\end{figure}
\noindent
Die relative Leistung besitzt ihr Maximum bei $r = 1$, was der Bedingung
$G_L = G_I$ entspricht. Damit wird erneut bestätigt, dass die maximale Leistungs-
übertragung bei Leistungsanpassung auftritt.

\noindent
Für sehr kleine Lastleitwerte ($r \ll 1$) sowie für sehr große Lastleitwerte
($r \gg 1$) geht die relative Leistung gegen null. Die Darstellung der relativen
Leistung erfolgt mit logarithmischer Skalierung der $r$-Achse, wodurch das Verhalten
der Leistung über mehrere Größenordnungen des Lastleitwerts übersichtlich dargestellt
werden kann.
